\section{Effective QCD with Auxiliary ``Heavy-Quark'' Field}

To be concrete, consider the following non-local operator
\begin{equation}\label{Eq:nonlocalop}
 O(x,y) = \overline{\psi}(x) \Gamma L(x,y)\psi(y),
\end{equation}
where $\psi, \bar\psi$ denote the bare quark fields and $\Gamma$ a Dirac structure. $L(x, y)$ is a path-ordered gauge link from $y$ to $x$,
\begin{equation}
             L(x, y) = {\cal P}\exp\left(-ig\int^1_0 d\lambda\,\frac{dz^\mu}{d\lambda}  A_\mu(z(\lambda))\right)
\end{equation}
carrying a $(3,\bar 3)$ representation in color SU(3) group, where the path is parametrized such that $z(0)=y$ and $z(1)=x$. The gauge link ensures gauge invariance of the operator $O(x,y)$.
For a straight-line gauge link which is of our main interest in the present paper, we can choose $z^\mu(\lambda)=y^\mu+\eta\lambda\, n^\mu$ with $n^\mu$ a unit vector specifying the direction and $\eta$ the length of the gauge link.

To study the renormalization property of the above operator, we introduce an ``heavy quark'' auxiliary field (color source without spin degrees of freedom)
$Q$ with color 3,
and its conjugate ${\overline Q}$, with color ${\bar 3}$.
We extend QCD to include this ``heavy quark'' interaction with the gluon field,  and introduce the following Lagrangian
\beq\label{Eq:lag}
{\cal L}={\cal L_{\text{QCD}}}+\overline Q(x)i n\cdot D Q(x),
\eeq
where $D_\mu=\partial_\mu+igt^a A^a_\mu$ is the covariant derivative in fundamental representation. For a real heavy quark or timelike Wilson line, $n^\mu$ is a timelike vector $n^\mu=(1,0,0,0)$, whereas for a spacelike Wilson line $n^\mu$ can be chosen as $n^\mu=(0,0,0,1)$. We will focus on the latter case in the following, although the discussion may go through for any $n$.
So long as there is a non-zero time component $n^0$, the ``heavy quark'' is dynamical, but follows a designated world line.

In the above theory, we can replace the bilocal operator $O(x,y)$ by a new operator,
\begin{equation}\label{Eq:repl}
             O(x, y) = \overline{\psi}(x)\Gamma Q(x) \overline{Q}(y) \psi(y).
\end{equation}
This can be seen as following. After integrating out the ``heavy quark'' field, we have
\beq
\int {\cal D}\overline Q{\cal D}Q\, Q(x)\overline Q(y) e^{i\int d^4x {\cal L}}=S_Q(x, y) e^{i\int d^4x {\cal L_{\text{QCD}}}}
\eeq
up to a determinant ${\rm det}(in\cdot D)$ which can be shown to be a constant and absorbed into the normalization of the generating functional~\cite{Mannel:1991mc}, where $S_Q(x, y)$ satisfies
\beq
n\cdot D\, S_Q(x, y)=\delta^{(4)}(x-y),
\eeq
with the solution in the case of $n^\mu=(0,0,0,1)$
\begin{align}
&S_Q(x,y)=\theta(x^z-y^z)\delta(x^0-y^0)\delta^{(2)}(\vec x_\perp-\vec y_\perp)L(x,y)\non\\
&=\theta(x^z-y^z)\delta(x^0-y^0)\delta^{(2)}(\vec x_\perp-\vec y_\perp)L(x^z,y^z),
\end{align}
under an appropriate boundary condition. In the following we will restrict to the case $x^z>y^z$ without loss of generality. The $\delta$-functions ensure that the time and transverse components of $x$ and $y$ are equal, and therefore generate a spacelike Wilson line along the longitudinal direction. The above result allows us to replace the Wilson line $L(x^z,y^z)$, which is the one appearing in the quasi-PDF, by the product of two auxiliary heavy quark fields $Q(x^z)\overline Q(y^z)$. The non-local operator in the quasi-PDF then becomes the product of two composite operators in Eq.~(\ref{Eq:repl}).
